%% CPP 2027 Submission — 0pat: The Prime Axis as a Circle (Fermat's Little Theorem)
\documentclass[sigplan,10pt,anonymous,review]{acmart}
\settopmatter{printfolios=true,printacmref=false}

\AtBeginDocument{%
  \providecommand\BibTeX{{\normalfont B\scshape ib\TeX}}}

\settopmatter{printacmref=false}
\renewcommand\footnotetextcopyrightpermission[1]{}
\pagestyle{plain}

\usepackage{microtype}
\usepackage{booktabs}
\usepackage{amsmath}
% XeTeX (tectonic): acmart loads newtxmath/libertine which already defines \Bbbk;
% reset it so amssymb's definition is accepted.
\let\Bbbk\relax
\usepackage{amssymb}

\begin{document}

\title{The Prime Axis as a Circle:\\ Fermat's Little Theorem via the
Multiplicative Group of $\mathbb{Z}/p\mathbb{Z}$ --- a Pure-Known-Mathematics
Re-formalization}

\author{Anonymous}
\affiliation{%
  \institution{Anonymous}
}
\email{anonymous@example.org}

\begin{abstract}
We report the second completed exercise of a re-formalization program
(``0pat''): a pat framework insight (``the prime axis, periodized, is a
circle; under the Euler-circle view the prime axis itself is a circle'') is
re-stated in pure known mathematics and formally proved in Lean 4 / mathlib
with \emph{no} pat framework concepts and \emph{no} complex plane. The
re-statement is the algebra of the multiplicative circle group
$(\mathbb{Z}/p\mathbb{Z})^\times$: the circle's size ($p-1$ elements),
the circle's algebra (Fermat's little theorem, both the coprime form
$a^{p-1} \equiv 1$ and the complete form $a^p \equiv a \pmod p$), and the
global structure (infinitely many primes). The proof is four theorems built
entirely on mathlib: an explicit bijection
$\{x \ne 0\} \simeq \mathrm{Fin}(p-1)$, Lagrange's theorem via
\texttt{orderOf\_dvd\_card} and \texttt{pow\_orderOf\_eq\_one}, the
congruence-lifting \texttt{Nat.ModEq.pow}, and
\texttt{Nat.exists\_infinite\_primes}. All four compile with zero errors
and zero \texttt{sorry} against mathlib v4.32.2.
\end{abstract}

\keywords{Lean, formalization, Fermat's little theorem, modular arithmetic, circle group, 0pat, reproducibility}

\maketitle

\section{Introduction}

\emph{Program.} This paper is the second completed exercise of a
re-formalization program (``0pat''): take the exercises of the pat
framework --- a body of Lean-formalized observations about basepoint-
relative operator semantics --- and re-prove their mathematical content
using \emph{pure known mathematics}: only mathlib's existing definitions
and theorems, no pat framework concepts (basepoints, phase locking,
interlock pairs, projection superpositions).

\emph{Method.} Each exercise is located on a genealogy map (rows: pat
concepts R1--R25; columns: mathematical genealogies). The map's entries
are the paths: a pat concept at coordinate $(R_i, C_j)$ has a
standard-mathematics correspondent at that coordinate. The
re-formalization walks the path: re-state the pat claim at its coordinate
in standard mathematics, formalize the equivalence, and prove.

\emph{This exercise.} The pat insight under study arose in the Riemann-
direction work: the number line, periodized modulo $p$, curls into the
circle $\mathbb{Z}/p\mathbb{Z}$; in this ``Euler-circle'' view the prime
axis itself is a circle, and its structure should be observable by the
same method used for the Riemann direction --- without needing the complex
plane. The insight sits at $(R23, C3)$ on the genealogy map (prime circle
$\times$ analytic number theory). Its standard-mathematics correspondent
is the algebra of the multiplicative circle group
$(\mathbb{Z}/p\mathbb{Z})^\times$: Fermat's little theorem and the
infinitude of primes. We re-state and prove it.

\section{Re-statement (0pat)}

\paragraph{Pat claim.} The prime axis, periodized, is a circle; on the
circle, the prime structure is the multiplicative circle group: every
nonzero residue is invertible (the circle is a group of size $p-1$), and
the circle's algebra is $x^p = x$ (Fermat's little theorem). The prime
axis is not exhausted by any finite modulus: there are infinitely many
primes.

\paragraph{0pat re-statement.} Let $p$ be prime.
\begin{enumerate}
\item \emph{Circle size:} $|(\mathbb{Z}/p\mathbb{Z})^\times| = p - 1$
  (the field $\mathbb{Z}/p\mathbb{Z}$ is a domain, so every nonzero
  element is a unit).
\item \emph{Circle algebra (coprime form):} if $\gcd(a,p) = 1$, then
  $a^{p-1} \equiv 1 \pmod p$ --- the exponent of the circle group divides
  the group's order (Lagrange).
\item \emph{Circle algebra (complete form):} for \emph{every} $a$,
  $a^p \equiv a \pmod p$ --- including $p \mid a$, where both sides
  vanish modulo $p$.
\item \emph{Global structure:} for every $n$, there is a prime $p \ge n$
  (the prime axis is infinite).
\end{enumerate}

\section{Formalization}

The proof is a single Lean file (\texttt{proof.lean}, 212 lines), importing
\texttt{Mathlib.Data.ZMod.Basic}, \texttt{Mathlib.GroupTheory.OrderOfElement},
and \texttt{Mathlib.Data.Nat.Prime.Basic}. Four theorems:

\begin{center}
\begin{tabular}{lll}
\toprule
Theorem & Statement & mathlib anchor \\
\midrule
\texttt{card\_units\_zmod\_prime} & $|(\mathbb{Z}/p\mathbb{Z})^\times| = p-1$
  & explicit bijection $\{x\ne0\} \simeq \mathrm{Fin}(p-1)$ \\
\texttt{fermat\_little} & $\gcd(a,p)=1 \Rightarrow a^{p-1} \equiv 1 \pmod p$
  & \texttt{orderOf\_dvd\_card}, \texttt{pow\_orderOf\_eq\_one} \\
\texttt{fermat\_little\_complete} & $\forall a,\ a^p \equiv a \pmod p$
  & \texttt{Nat.ModEq.pow}, \texttt{mod\_eq\_zero\_of\_dvd} \\
\texttt{infinitely\_many\_primes} & $\forall n,\ \exists p \ge n$, prime $p$
  & \texttt{Nat.exists\_infinite\_primes} \\
\bottomrule
\end{tabular}
\end{center}

\paragraph{Circle size by explicit bijection.} Instead of appealing to
mathlib's unit-coprime equivalence, the proof constructs the bijection
$\{x : \mathbb{Z}/p\mathbb{Z} \mid x \ne 0\} \simeq \mathrm{Fin}(p-1)$
directly: $x \mapsto x.\mathrm{val} - 1$ (residues $1 \ldots p-1$ to
indices $0 \ldots p-2$), inverse $k \mapsto (k+1)$. Both directions are
checked; the range proof uses that a nonzero residue has value $0 <
x.\mathrm{val} < p$, so $x.\mathrm{val} - 1 < p - 1$. This makes the
``circle of size $p-1$'' statement self-contained.

\paragraph{Fermat via Lagrange.} A coprime $a$ maps to the unit group via
\texttt{unitOfCoprime}; the order of that unit divides the group order by
Lagrange (\texttt{orderOf\_dvd\_card}), and
\texttt{pow\_orderOf\_eq\_one} gives $u^{k} = 1$ for $k =
\mathrm{order}(u) \mid p-1$. The congruence is recovered by casting back
from $(\mathbb{Z}/p\mathbb{Z})^\times$ to $\mathbb{Z}/p\mathbb{Z}$ and
from $\mathbb{Z}/p\mathbb{Z}$ to $\mathbb{N}$ via
\texttt{natCast\_eq\_natCast\_iff}.

\paragraph{Complete form by congruence lifting.} The non-coprime case
($p \mid a$) is handled by the observation that $a \equiv 0 \pmod p$
(\texttt{Nat.mod\_eq\_zero\_of\_dvd}), lifted to powers by
\texttt{Nat.ModEq.pow}: $a^p \equiv 0^p = 0 \pmod p$. Both sides
congruent to $0$ gives $a^p \equiv a$. This is the cleanest expression of
``on the circle, $p \mid a$ means the point is the zero of the ring''.

\paragraph{Verification.} Compiled with Lean 4.32.2 (mathlib v4.32.2),
zero errors, zero \texttt{sorry}, zero axioms beyond mathlib's. Build
command in the artifact.

\section{Relation to the pat exercise}

The original insight (pat exercise XXIV) was formulated inside the pat
framework. This re-formalization shows that the mathematical content
survives the removal of the framework --- and of the complex plane: the
``prime axis as a circle'' is exactly the modular circle
$\mathbb{Z}/p\mathbb{Z}$ and its multiplicative group, entirely
first-order algebra over a finite field. The pat vocabulary is preserved
in the exercise record (README) as provenance, not in the proof.

\section{Conclusion}

The second 0pat exercise is complete: the prime axis as a circle ---
Fermat's little theorem, the size of the multiplicative circle group, and
the infinitude of primes --- proved with pure known mathematics in four
theorems, zero errors. The observation requires no complex analysis: the
circle structure of the primes modulo $p$ is the algebra of
$(\mathbb{Z}/p\mathbb{Z})^\times$. The method (genealogy-map
path-walking) is validated on a second, algebraically different exercise.

\bibliographystyle{ACM-Reference-Format}
\bibliography{refs}

\end{document}
